
\section{S.M.A.R.T. Goals, Todo}

Todo
Work for future Roland


\section{S.M.A.R.T. Goals , TODO}
\label{ResearchSmartGoals}

\begin{itemize}





\item \textbf{1) Suggested Goals}
1. Definition of the Research and Evaluation Framework
PHASE 1 — RESEARCH

WHY: To establish a scientifically grounded and reproducible approach for evaluating automated data quality and availability checks.

By the end of the research phase, the selected use case, research question, relevant validation requirements and evaluation methodology are defined. This includes the criteria required to compare manual and automated data quality validation approaches.

Success is measured by a documented research and evaluation framework specifying the selected data quality metrics, comparison methods, required data and evaluation criteria.



\item \textbf{2) Suggested Goals}

2. Establishment of a Baseline for Data Quality Validation
PHASE 2 — DEFINE THE BASELINE WITH STAKEHOLDERS

WHY: To establish a reference against which the effectiveness of automated data quality validation can be evaluated.

By the end of the baseline phase, the current manual data quality validation process is documented, and a systematic approach for recording identified data quality issues is established with business stakeholders. A defined observation period is prepared for comparing manual and automated validation results.

Success is measured by a documented baseline process, agreed error recording criteria, an established recording mechanism and a defined comparison period.

Passender Meilenstein: M3, 29.11.2026 – Error documentation started.

\item \textbf{3) Suggested Goals}
3. Implementation and Technical Validation of Automated DQM Checks
PHASE 3 — FIRST ITERATION

WHY: To demonstrate that the selected reporting data can be automatically validated using the existing technical foundation and the defined business rules.

By the end of the first implementation phase, an automated data quality validation solution for the selected use case is implemented and tested. The solution must execute the defined validation rules on representative reporting data and record the resulting validation outcomes.

Success is measured by a reproducible implementation, successful execution of the selected validation checks across multiple test runs, and documented technical validation results.

Passender Meilenstein: M6, 21.03.2027 – MVP established; M7, 04.04.2027 – stakeholder presentation.

\item \textbf{4) Suggested Goals}
4. Development and Validation of an Improved or Alternative Approach
PHASE 4 — SECOND ITERATION

WHY: To investigate whether limitations identified during the first iteration can be addressed through an improved or alternative validation approach.

By the end of the second implementation phase, the results and limitations of the first iteration are analysed, and at least one improvement or alternative approach is implemented and tested using comparable validation criteria.

Success is measured by a documented analysis of the first iteration, a justified selection of the second approach and reproducible test results for both iterations.

Passender Meilenstein: M8, 18.04.2027 – second iteration defined. Der Abschluss der Implementierung muss vor dem Vergleich in Phase 5 erfolgen.

\item \textbf{5) Suggested Goals}
5. Comparative Evaluation of Automated Data Quality Validation
PHASES 5–6 — COMPARATIVE EVALUATION AND FINALISATION

WHY: To determine the effectiveness and limitations of the investigated approaches and provide an evidence-based answer to the research question.

By the end of the comparative evaluation phase, the investigated validation approaches are systematically compared against the defined baseline and evaluation criteria. The results are analysed to determine their detection capability and operational applicability within the selected reporting use case.

Success is measured by a documented comparison of validation results, quantitative measurements for the selected evaluation criteria, identification of limitations and a scientifically justified assessment of the research question.

Passender Meilenstein: M9, 16.05.2027 – approaches compared. Die Interpretation und Schlussfolgerungen werden bis M10, 11.07.2027, abgeschlossen.

\end{itemize}